В ходе выпускной квалификационной работы была достигнута поставленная цель:
разработаны CLI-инструмент и веб-интерфейс для системы нагрузочного тестирования
HTTP и gRPC сервисов, обеспечивающие гибкое управление нагрузкой
и визуализацию метрик в реальном времени.

В процессе выполнения работы были решены следующие задачи.

Проведён сравнительный анализ существующих инструментов нагрузочного
тестирования~--- Apache JMeter, Gatling, k6, Vegeta и hey.
В ходе анализа установлено, что ни один из рассмотренных инструментов
не сочетает в себе распространение в виде единого бинарного файла, поддержку HTTP и gRPC
в одном дистрибутиве и встроенный веб-интерфейс с потоковой передачей метрик,
что обусловило необходимость разработки инструмента Tsunami.

Сформулированы функциональные и нефункциональные
требования к разрабатываемому инструменту.
Среди ключевых нефункциональных требований: поставка единым бинарным файлом,
кросс-платформенная поддержка шести платформ, целевая производительность
не менее 1000~RPS при 50 параллельных исполнителях и задержка WebSocket-трансляции
не более 100~мс.

Спроектирована трёхслойная архитектура системы, включающая движок нагрузочного
тестирования, CLI-слой на базе Cobra и HTTP-сервер с WebSocket-сервером.
Для формализации проектных решений разработаны диаграммы прецедентов, классов,
состояний и последовательности.

Реализован движок нагрузочного тестирования с поддержкой HTTP и gRPC.
Движок построен по паттерну «резерв исполнителей» с ограничителем скорости
на алгоритме текущего ведра.
Потокобезопасный сбор метрик обеспечивается кольцевым буфером ёмкостью
100\,000~замеров; по завершении теста вычисляются перцентили задержки
P50, P90, P95 и P99 за одну сортировки.
Реализована классификация сетевых ошибок по пяти категориям,
что позволяет оперативно отличить инфраструктурные сбои от прикладных ошибок сервера.

Реализованы CLI-интерфейс и веб-интерфейс системы для управления нагрузочными
тестами и мониторинга метрик в реальном времени.
CLI-интерфейс на базе Cobra включает команды \texttt{attack}, \texttt{serve}
и \texttt{grpc} с настраиваемой интенсивностью запросов и штатным завершением теста.
Веб-интерфейс реализован как React SPA с графиками Recharts и встраивается
в исполняемый файл через директиву \texttt{//go:embed};
REST~API из восьми маршрутов и WebSocket-сервер с интервалом трансляции 50~мс
обеспечивают управление тестами и потоковую передачу метрик через браузер.

Проведено двухуровневое тестирование системы.
Юнит-тесты функции \texttt{CalculateAllPercentiles} охватывают 13~сценариев,
включая граничные случаи.
Интеграционные тесты с библиотекой testcontainers-go используют реальный
Docker-контейнер с сервисом httpbin и верифицируют точность классификации ошибок,
корректность счётчиков метрик и соответствие REST~API заявленным контрактам.
Все тесты успешно проходят в среде GitHub Actions~CI.

\textbf{Практическая ценность} разработанного инструмента Tsunami состоит в том,
что он обеспечивает нагрузочное тестирование HTTP и gRPC сервисов в виде
единого самодостаточного бинарного файла без зависимостей среды выполнения,
распространяемого для шести целевых платформ через GitHub Releases и Homebrew.
Инструмент готов к применению инженерами по качеству и DevOps-специалистами
в профессиональной деятельности, в том числе в составе автоматизированных
конвейеров CI/CD.

Перспективными направлениями развития инструмента являются поддержка сценарного
нагрузочного тестирования с переменной интенсивностью, интеграция трассировки
через OpenTelemetry и реализация распределённого режима нагрузки
с несколькими агентами.
